\documentclass[12pt]{article}
\usepackage[utf8]{inputenc}
\usepackage[spanish]{babel}
\usepackage{amsmath}
\usepackage{amssymb}
\usepackage{braket} % Paquete ideal para la notación de Dirac
\usepackage{geometry}
\geometry{a4paper, margin=2.5cm}

\title{Notas sobre la Caminata Cuántica Discreta}
\author{Saúl Nájera Allara \\ \small Escuela de Ciencias Físicas y Matemáticas, USAC}
\date{\today}

\begin{document}

\maketitle

\section*{Evolución del Caminante Cuántico}

Sea $\mathbb{H_p}$ el espacio de hilbert de la posición y $\mathcal{H_c}$ el espacio de hilbert de la moneda. El estado del caminante 
pertenece al espacio:
\begin{equation}
\mathcal{H}=\mathcal{H}_p \otimes \mathcal{H}_c
\end{equation}
Considerando que $dim\, \mathcal{H}_c=4$
\end{document}